\section{Problema detectado y/o faltante}

En el panorama actual no existe una solución que integre, de forma accesible y robusta, la creación de mundos al estilo RPG Maker con la infraestructura masiva y persistente de un MMO. 
Las plataformas existentes presentan limitaciones: motores como Unity o RPG Maker no contemplan de manera nativa escenarios multijugador de gran escala, 
mientras que títulos como Roblox o Minecraft, si bien ofrecen creación comunitaria, imponen barreras técnicas o carecen de un enfoque de rol persistente.  
\vspace{0.25cm}

Los principales problemas a resolver incluyen: \textbf{desarrollar un editor intuitivo y accesible} para usuarios sin experiencia en programación, 
pero lo suficientemente potente para permitir personalización avanzada. En este entorno, los creadores deben poder diseñar ítems, \textit{quests},
NPCs con comportamientos interactivos y enemigos con dinámicas de combate, asegurando variedad y evitando experiencias 
monótonas o repetitivas.; \textbf{definir incentivos} claros que motiven tanto a creadores como a jugadores a participar activamente; 
diseñar un \textbf{sistema de monetización} sostenible que no desincentive el uso, ni genere desequilibrios económicos; 
y establecer una arquitectura de servidores escalable, tolerante a fallos y de baja latencia, capaz de sostener múltiples mundos simultáneos con interacción en tiempo real.  
\vspace{0.25cm}

En síntesis, el faltante identificado es la ausencia de un ecosistema que combine facilidad de creación, incentivos adecuados, monetización justa e infraestructura sólida. 
El proyecto busca abordar estas brechas a través de un MVP que permita validar las bases técnicas y de diseño necesarias para su evolución futura.
